\section{Propagation of low-degree perturbations}
\label{sec:propagation}

Pauli expansions and low-degree approximations have been used to study
$\QAC$ \citep{NadimpalliParhamVasconcelosYuen2024,
AnshuDongOuYao2025,DongOuYao2025}.
Here regularity lets us replace the large product projections arising
from conjugation by polynomials of small degree.  The resulting
approximation controls a perturbed vector along an exact reference
trajectory.

\begin{proposition}[Low-degree propagation]
\label{prop:propagation}
Let $V$ be a circuit with a bounded number of analytic layers, and suppose
the trajectory from $\zeta$, with $\norm{\zeta}\le 1$, is regular for $V$.
Here $k,K$ denote degree exponents: fix $0<k<K<1$.  For every operator $E$
with
\[
 \deg E\le m^k,
 \qquad \norm{E}\le \exp(m^k),
\]
there is, eventually, an operator $E'$ such that
\begin{equation}
 \label{eq:propagation}
 \begin{gathered}
  \norm{VE\zeta-E'V\zeta}\le \exp(-m^k),\\
  \deg E'\le m^K,
  \qquad \norm{E'}\le \exp(m^K).
 \end{gathered}
\end{equation}
The starting vector need not be a product vector.
\end{proposition}

\begin{proof}
We first prove the one-layer assertion with the stronger error
$o(\exp(-m^k))$.  Conjugation by a layer of one-qubit unitaries preserves
both degree and operator norm and gives zero error.  It remains to consider
a layer
\[
 F=\prod_g (I-2P_g)
\]
of reflections on pairwise disjoint supports.  Set $\xi=F\zeta$.

\emph{A controlled local expansion.}
On one qubit use the unitary matrices $I,X,Z,XZ$, where
\[
 X=\begin{pmatrix}0&1\\1&0\end{pmatrix},
 \qquad
 Z=\begin{pmatrix}1&0\\0&-1\end{pmatrix}.
\]
Their tensor products form an orthonormal basis for the inner product
$2^{-m}\Tr(A^\dagger B)$.  Thus the unique expansion of $E$ is
\[
 E=\sum_D c_DD,
 \qquad c_D=2^{-m}\Tr(D^\dagger E).
\]
Put $\ell=\lfloor m^k\rfloor$.  Every summand in any degree-$\ell$
decomposition of $E$ has zero coefficient against a tensor $D$ whose
support is not contained in that summand's support: a nonidentity factor
of $D$ outside the support has trace zero.  Consequently $c_D=0$ whenever
$\abs{\supp D}>\ell$.  This conclusion requires no bound on the norms or
coefficients of the original decomposition.  Also,
\[
 \abs{c_D}
 \le 2^{-m}\,2^m\norm{D^\dagger E}
 =\norm{E},
\]
because $\abs{\Tr M}\le 2^m\norm{M}$ for every operator on $m$ qubits.
There are at most
\[
 N_m=\sum_{j=0}^{\ell}\binom mj3^j
 \le (\ell+1)(3m)^\ell
\]
possible nonzero coefficients, and hence
$\sum_D\abs{c_D}\le \exp(m^k)N_m$.

Fix one such $D$, supported on $J$, and let $\mathcal T$ be the gates
whose supports meet $J$.  Disjointness gives
$\abs{\mathcal T}\le\abs J\le\ell$.  All other reflections commute with
$D$ and cancel in $FDF^\dagger$.  With
$P_{\mathcal L}=\prod_{g\in\mathcal L}P_g$ and
empty products interpreted as identities, expansion on both sides gives
\[
 FDF^\dagger
 =\sum_{\mathcal L,\mathcal R\subseteq\mathcal T}
   (-2)^{\abs{\mathcal L}+\abs{\mathcal R}}
   P_{\mathcal L}DP_{\mathcal R}.
\]
Factor the Hilbert space across $J$ and $J^c$, writing
$D=D_J\otimes I$ and $P_g=A_g\otimes B_g$.
The factors $A_g$ and $B_g$ are the corresponding product projections,
possibly identities.  If
$A_{\mathcal L}=\prod_{g\in\mathcal L}A_g$, then each term factors as
\[
 P_{\mathcal L}DP_{\mathcal R}=M\otimes Q,
 \qquad
 M=A_{\mathcal L}D_JA_{\mathcal R},
 \qquad
 Q=\prod_{g\in\mathcal L\cup\mathcal R}B_g.
\]
Indeed, the exterior factors have disjoint supports for distinct gates,
and repeated factors satisfy $B_g^2=B_g$.  The operator
$M$ has norm at most one.  Its displayed order is
retained; no projection on $J$ has been commuted through $D_J$.
The sum of the absolute scalar coefficients in this expansion is
\[
 \sum_{\mathcal L,\mathcal R\subseteq\mathcal T}
  2^{\abs{\mathcal L}+\abs{\mathcal R}}
 =9^{\abs{\mathcal T}}\le 9^\ell.
\]

\emph{Replacing the exterior projections.}
Let $H_Q$ count failures of the one-qubit factors of $Q$,
extended by identity outside their sites.  Thus
$Q=\one_{H_Q=0}$ on $J^c$; when this test is empty its count is
zero.  The operator $H_Q$ is a partial sum of the commuting site
counts in $H_{\mathcal T}$.  In their common eigenbasis, its value is at
most that of $H_{\mathcal T}$.

Choose fixed exponents
\[
 k<p_1<L<K,
 \qquad
 \frac{L+p_1}{2}<r<p_2<L.
\]
Such a choice exists since $(L+p_1)/2<L$; in particular,
$k<p_1<r<p_2<L<K$.  Regularity, used with
$(b,h,p)=(k,L,p_2)$, holds uniformly over $\mathcal T$ and gives
\[
 \norm{\one_{H_Q>m^L}\xi}
 \le \norm{\one_{H_{\mathcal T}\ge m^L}\xi}
 \le \exp(-m^{p_2}).
\]
Here the reference vector is exactly $F\zeta$, before applying
$M$.

By Lemma~\ref{lem:zero-filter}, choose a real polynomial $f$ of degree
at most $m^r$, bounded in absolute value by $\exp(m^r)$ on $[0,m]$,
whose error in approximating $\one_{x=0}$ is at most
$\exp(-m^{p_1})$ at every integer $0\le x\le m^L$.
The spectrum of $H_Q$ is integral.  Splitting $\xi$ into counts at
most $m^L$ and larger counts therefore yields
\begin{align*}
 \norm{(Q-f(H_Q))\xi}
 &\le \exp(-m^{p_1})\norm{\xi}
    +\bigl(1+\exp(m^r)\bigr)
       \norm{\one_{H_Q>m^L}\xi}\\
 &\le \exp(-m^{p_1})
    +\bigl(1+\exp(m^r)\bigr)\exp(-m^{p_2}).
\end{align*}
In this display $Q$ and $f(H_Q)$ act as identity on $J$.
Applying the contraction $M$ afterwards cannot increase the
error.

Define $E_F$ by replacing each exterior $Q$ in the expansion
of $FEF^\dagger$ by $f(H_Q)$, retaining its coefficient and its
factor $M$.  The total absolute coefficient mass is at most
\[
 \Lambda_m
 =\exp(m^k)(\ell+1)(27m)^\ell,
 \qquad \log\Lambda_m=O(m^k\log m).
\]
It follows that
\begin{align*}
 \norm{FE\zeta-E_FF\zeta}
 &\le \Lambda_m\left[
       \exp(-m^{p_1})
       +\bigl(1+\exp(m^r)\bigr)\exp(-m^{p_2})
       \right]\\
 &\le 3\exp\bigl(-\tfrac12m^{p_1}\bigr)
   =o\bigl(\exp(-m^k)\bigr)
\end{align*}
eventually.  To check the second inequality, use
$\log\Lambda_m=o(m^{p_1})$, $r<p_2$, and $p_1<p_2$;
each of the two bracketed contributions, after multiplication by
$\Lambda_m$, is absorbed by the displayed bound.

Each new term has degree at most $\ell+m^r$: its factor on $J$ has
support size at most $\ell$, and expanding the polynomial in one-site
counts gives degree at most $m^r$ on the remaining sites.
Its norm is at most $\exp(m^r)$.  Thus, eventually,
\[
 \deg E_F\le\ell+m^r\le m^K,
 \qquad
 \norm{E_F}\le\Lambda_m\exp(m^r)\le\exp(m^K),
\]
because $k<r<K$ and $\log\Lambda_m=O(m^k\log m)$.
This proves the stronger one-layer assertion, uniformly over the
admissible $E$.

\emph{A bounded number of layers.}
Pad by identity layers so that $V=U_t\cdots U_1$ has a fixed number
$t$ of layers.  The case $t=0$ is immediate with $E'=E$.
For $t\ge1$, choose a fixed increasing chain
\[
 k=k_0<k_1<\cdots<k_t=K.
\]
Let $\zeta_0=\zeta$ and $\zeta_i=U_i\zeta_{i-1}$ be the exact regular
reference trajectory.  Set $E_0=E$ and apply the one-layer result
successively with exponents $(k_{i-1},k_i)$ to obtain $E_i$ with
degree at most $m^{k_i}$, norm at most $\exp(m^{k_i})$, and local
residual
\[
 \delta_i
 =U_iE_{i-1}\zeta_{i-1}-E_i\zeta_i,
 \qquad
 \norm{\delta_i}=o\bigl(\exp(-m^{k_{i-1}})\bigr).
\]
For the exact perturbed evolution $\eta_0=E\zeta$,
$\eta_i=U_i\eta_{i-1}$, put $r_i=\eta_i-E_i\zeta_i$.
Then $r_0=0$ and the exact recurrence is
\[
 r_i=U_i r_{i-1}+\delta_i.
\]
Unitarity implies
\[
 \norm{r_t}\le\sum_{i=1}^t\norm{\delta_i}
 =o\bigl(\exp(-m^k)\bigr)
 \le\exp(-m^k)
\]
eventually.  In particular, no later approximating operator acts on an
earlier residual.  Taking $E'=E_t$ proves
\eqref{eq:propagation}.
\end{proof}
